\documentclass[11pt,a4paper]{article}

\usepackage[utf8]{inputenc}
\usepackage[T1]{fontenc}
\usepackage[spanish]{babel}
\usepackage{amsmath,amssymb,amsthm,bm}
\usepackage{mathtools}
\usepackage{geometry}
\usepackage{booktabs}
\usepackage{array}
\usepackage{tabularx}
\usepackage{float}
\usepackage{enumitem}
\usepackage{xcolor}
\usepackage{tcolorbox}
\usepackage{hyperref}
\usepackage{caption}
\usepackage{tikz}
\usepackage{pgfplots}
\pgfplotsset{compat=1.18}
\usetikzlibrary{arrows.meta,positioning,shapes.geometric,calc,fit,backgrounds,decorations.pathreplacing,patterns}

\geometry{margin=2.5cm}

\definecolor{cblue}{RGB}{41,128,185}
\definecolor{cgreen}{RGB}{39,174,96}
\definecolor{cred}{RGB}{231,76,60}
\definecolor{corange}{RGB}{243,156,18}
\definecolor{cpurple}{RGB}{142,68,173}
\definecolor{cgray}{RGB}{149,165,166}
\definecolor{cdark}{RGB}{44,62,80}
\definecolor{clight}{RGB}{236,240,241}

\tcbset{
  defbox/.style={colback=cblue!5, colframe=cblue!60, fonttitle=\bfseries, arc=2mm, boxrule=0.5pt},
  thmbox/.style={colback=cgreen!5, colframe=cgreen!60, fonttitle=\bfseries, arc=2mm, boxrule=0.5pt},
  warnbox/.style={colback=corange!5, colframe=corange!60, fonttitle=\bfseries, arc=2mm, boxrule=0.5pt},
  keybox/.style={colback=cpurple!5, colframe=cpurple!60, fonttitle=\bfseries, arc=2mm, boxrule=0.5pt},
  redbox/.style={colback=cred!5, colframe=cred!60, fonttitle=\bfseries, arc=2mm, boxrule=0.5pt},
}

\newtheorem{definition}{Definici\'on}[section]
\newtheorem{theorem}{Teorema}[section]
\newtheorem{lemma}{Lema}[section]
\newtheorem{corollary}{Corolario}[section]
\newtheorem{proposition}{Proposici\'on}[section]
\newtheorem{remark}{Observaci\'on}[section]

\newcommand{\E}{\mathbb{E}}
\newcommand{\Prob}{\mathbb{P}}
\newcommand{\R}{\mathbb{R}}
\newcommand{\KL}{D_{\mathrm{KL}}}
\newcommand{\af}{\alpha_f}
\newcommand{\score}{\hat{S}_{\af}}
\newcommand{\energy}{E_{\af}}
\newcommand{\exh}{\hat{E}}
\newcommand{\param}[1]{\texttt{\small\color{cpurple}#1}}

\begin{document}

\begin{center}
    {\LARGE\bfseries Teor\'ia Unificada del Sistema}\\[4pt]
    {\LARGE\bfseries $\af$ Bifurcation Short + AEPS}\\[6pt]
    {\large Modelo de Mezcla, Super-Aditividad,}\\[2pt]
    {\large Convergencia y Monitoreo de Salud}\\[16pt]
    {\large Javier Epeloa\quad$\cdot$\quad Mapplics}\\[8pt]
    {\normalsize Documento T\'ecnico --- Marzo 2026}\\[4pt]
    {\small\itshape Formalizaci\'on matem\'atica completa del sistema de trading:\\
    desde la bifurcaci\'on f\'isica hasta las garant\'ias de convergencia}
\end{center}

\vspace{8pt}

\begin{abstract}
Se presenta la formalizaci\'on matem\'atica completa del sistema
de trading $\af$ Bifurcation Short, integrando la se\~nal de entrada
con el sistema adaptativo de salida AEPS. Se demuestra que: (1) toda
se\~nal de entrada en mercados financieros genera inherentemente una
distribuci\'on bimodal de resultados, lo cual explica por qu\'e los
enfoques cl\'asicos de machine learning fracasan al intentar predecir
retornos; (2) el sistema AEPS explota esta bimodalidad mediante
clasificaci\'on secuencial post-entrada, reformulando el problema de
predicci\'on (irresoluble) como un problema de clasificaci\'on temporal
(resoluble); (3) la combinaci\'on se\~nal + AEPS es super-aditiva;
(4) se establecen formalmente las \textbf{tres condiciones} que
cualquier se\~nal de entrada debe cumplir para que AEPS genere edge;
(5) la convergencia del sistema se garantiza bajo condiciones
espec\'ificas y puede acelerarse mediante dry run pre-operativo;
y (6) un monitor de salud de tres alarmas detecta p\'erdida de
bimodalidad antes de da\~no significativo al capital. Se valida
emp\'iricamente sobre $N=46$ trades reales ($p = 0.000315$,
Mann-Whitney), demostrando profit factor 3.53 post-convergencia.
\end{abstract}

\tableofcontents
\newpage

% ═══════════════════════════════════════════════════════════════
\section{Fundamentos: La Bifurcaci\'on como Generador de Mezcla}
\label{sec:mixture}
% ═══════════════════════════════════════════════════════════════

\subsection{La Ecuaci\'on de Langevin del Mercado}

El sistema modela el precio de futuros perpetuos como una part\'icula
en un potencial con fuerza externa:

\begin{equation}
L\ddot{x} + \gamma\dot{x} + \kappa\,x = F_{\text{ext}} + \sigma\xi
\label{eq:langevin}
\end{equation}

donde $\kappa$ es la constante del spring (derivada del funding rate),
$F_{\text{ext}}$ es la fuerza externa (compra retail, leverage),
y $\sigma\xi$ es el ruido estoc\'astico. La se\~nal $\score$ detecta
la condici\'on $\af < 0$ (spring invertido), que indica que el
potencial est\'a en un m\'aximo local y el sistema puede transitar
a un nuevo equilibrio.

\subsection{El Modelo de Mezcla}

\begin{definition}[Modelo de Mezcla $\af$]
\label{def:mixture}
Sea $\score \geq \theta_{\min}$ la condici\'on de entrada. El vector
de trayectoria $\bm{X} = (\text{MFE}, |\text{MAE}|, t_{\text{hold}})$
observado al cierre sigue una distribuci\'on de mezcla:
\begin{equation}
p(\bm{X} \mid \score \geq \theta_{\min}) =
  \pi_W \cdot f_W(\bm{X}) + (1 - \pi_W) \cdot f_L(\bm{X})
\label{eq:mixture}
\end{equation}
donde $\pi_W = \Prob(\text{winner} \mid \text{entry})$ es la
probabilidad de bifurcaci\'on exitosa, $f_W$ y $f_L$ son las
densidades de las poblaciones winner y loser respectivamente.
\end{definition}

\begin{definition}[Interpretaci\'on F\'isica de las Poblaciones]
Las dos poblaciones corresponden a los dos desenlaces posibles
cerca del punto cr\'itico $\af = 0$:
\begin{itemize}[leftmargin=1.5em]
\item \textbf{Cuenca $W$}: $\kappa$ invierte signo, $F_{\text{ext}} \to 0$
  (exhaustion). La fuerza restauradora empuja a favor del short.
  El precio cae r\'apido (MFE alto, MAE bajo).
\item \textbf{Cuenca $L$}: $F_{\text{ext}}$ sigue dominando, $\kappa$
  no cambia de signo. El pump contin\'ua. El precio sigue subiendo
  (MFE bajo, MAE alto).
\end{itemize}
La incertidumbre sobre cu\'al cuenca dominar\'a es intr\'inseca a la
din\'amica no lineal cerca del punto cr\'itico.
\end{definition}


\subsection{Separabilidad de la Se\~nal}

\begin{definition}[Divergencia de Kullback-Leibler]
\label{def:kl}
La separabilidad de la se\~nal se mide por la divergencia KL
\cite{cover2006elements} entre las distribuciones marginales de MFE:
\begin{equation}
\Delta = \KL\!\left(f_W^{\text{MFE}} \,\|\, f_L^{\text{MFE}}\right)
= \int f_W^{\text{MFE}}(m) \log \frac{f_W^{\text{MFE}}(m)}{f_L^{\text{MFE}}(m)} \, dm
\label{eq:kl}
\end{equation}
$\Delta = 0$ cuando las distribuciones son id\'enticas (entrada random);
$\Delta > 0$ cuando son distinguibles.
\end{definition}

La divergencia KL mide cu\'anta \textit{informaci\'on} se pierde al
describir datos de $f_W$ usando $f_L$. Intuitivamente: si un trade
tiene MFE de 4\%, $f_W$ le asigna alta probabilidad y $f_L$ baja
probabilidad. El ratio $f_W(m)/f_L(m)$ es grande, y $\log$ de ese
ratio cuantifica la ``evidencia'' de que el trade es winner. El promedio
ponderado de esta evidencia sobre toda la distribuci\'on $f_W$ es $\Delta$.

Propiedades relevantes:
\begin{itemize}[leftmargin=1.5em]
\item $\Delta \geq 0$ siempre (desigualdad de Gibbs), con igualdad
  sii $f_W = f_L$ c.s.
\item $\Delta$ no es sim\'etrica:
  $\KL(f_W \| f_L) \neq \KL(f_L \| f_W)$.
  La primera mide el costo de creer $L$ cuando es $W$
  (false positive del abort); la segunda mide el costo inverso.
\item Valores de referencia: $\Delta < 0.1$ (indistinguibles),
  $0.1$--$0.3$ (separaci\'on d\'ebil), $0.3$--$0.7$ (separaci\'on
  moderada), $> 1.0$ (separaci\'on fuerte).
\end{itemize}


\subsection{Validaci\'on Emp\'irica ($N=46$ trades)}

\begin{table}[H]
\centering
\caption{Par\'ametros del modelo de mezcla sobre el dataset completo.}
\label{tab:mixture_full}
\begin{tabular}{lccccc}
\toprule
\textbf{Poblaci\'on} & $\bm{n}$ & $\overline{\textbf{MFE}}$ &
  $\overline{|\textbf{MAE}|}$ & \textbf{E-ratio} & $\overline{\textbf{PnL}}$ \\
\midrule
Winners ($f_W$) & 27 & 3.74\% & 0.93\% & \textbf{4.01} & $+1.97\%$ \\
Losers ($f_L$) & 19 & 1.81\% & 3.80\% & \textbf{0.48} & $-2.93\%$ \\
\midrule
Mezcla total & 46 & 2.94\% & 2.12\% & 1.40 & $-0.06\%$ \\
\bottomrule
\end{tabular}
\end{table}

\begin{table}[H]
\centering
\caption{Tests de separabilidad ($N=46$).}
\label{tab:sep_tests}
\begin{tabular}{lccl}
\toprule
\textbf{Test} & \textbf{Estad\'istico} & $\bm{p}$\textbf{-value} & \textbf{Significancia} \\
\midrule
Mann-Whitney $U$ (MFE) & $U = 418$ & $\bm{0.000315}$ & S\'i (0.03\%) \\
$\KL(f_W^{\text{MFE}} \| f_L^{\text{MFE}})$ & 0.250 nats & --- & Moderada \\
Binomial (MFE $\geq 1\%$, $p_0=0.5$) & 32/46 & $\bm{0.006}$ & S\'i (0.6\%) \\
\bottomrule
\end{tabular}
\end{table}

\begin{tcolorbox}[thmbox, title=Resultado 1: Bimodalidad Confirmada ($N=46$)]
La se\~nal $\score$ produce dos poblaciones estad\'isticamente
separables con alta significancia ($p = 0.000315$). La separabilidad
$\Delta = 0.250$ nats es estable entre la primera mitad ($\Delta = 0.405$)
y la segunda ($\Delta = 0.551$) del dataset. La se\~nal funciona
desde el trade 1 --- la bimodalidad es una propiedad de la entrada,
no del sistema de salida.
\end{tcolorbox}


\subsection{Informaci\'on Mutua y Contenido Informacional}

\begin{definition}[Informaci\'on Mutua]
Sea $O \in \{W, L\}$ el outcome y $M = \mathbb{1}[\text{MFE} \geq \tau]$
un indicador binario. La informaci\'on mutua \cite{cover2006elements}
mide cu\'anto reduce $M$ la incertidumbre sobre $O$:
\begin{equation}
I(M; O) = H(O) - H(O \mid M)
\end{equation}
\end{definition}

Con $\tau = 1\%$ y los datos actuales ($N=25$ en ventana AEPS):
\begin{align}
H(O) &= 0.989 \text{ bits} \\
\Prob(W \mid \text{MFE} \geq 1\%) &= 0.750 \\
\Prob(W \mid \text{MFE} < 1\%) &= 0.222 \\
I(M; O) &= 0.195 \text{ bits} \quad (19.7\% \text{ de la incertidumbre total})
\end{align}

Observar si el MFE cruza 1\% proporciona casi una quinta parte de
la informaci\'on necesaria para clasificar el trade. Este es el
fundamento de por qu\'e AEPS usa MFE como estad\'istico principal
de decisi\'on.


% ═══════════════════════════════════════════════════════════════
\section{Por Qu\'e Fracasa el Machine Learning Predictivo}
\label{sec:ml_failure}
% ═══════════════════════════════════════════════════════════════

\subsection{El Problema de la Bimodalidad Inherente}

Proponemos que \textbf{toda se\~nal de entrada en mercados financieros
genera inherentemente una distribuci\'on bimodal de resultados}. Esta
observaci\'on, consistente con los modelos de Markov regime-switching
de Hamilton \cite{hamilton1989regime}, explica por qu\'e los enfoques
cl\'asicos de ML fracasan sistem\'aticamente al intentar predecir
retornos en trading.

\begin{proposition}[Irreducibilidad de la Bimodalidad]
\label{prop:bimodal_irreducible}
Sea $\mathcal{S}$ cualquier se\~nal de entrada basada en observables
de mercado. Para los mismos valores de features $\bm{x}$, el retorno
$r$ sigue una distribuci\'on de mezcla:
\begin{equation}
p(r \mid \bm{x}) = \pi_W(\bm{x}) \cdot f_W(r) + (1 - \pi_W(\bm{x})) \cdot f_L(r)
\label{eq:bimodal_inherent}
\end{equation}
donde $f_W$ y $f_L$ son distribuciones con medias de signo opuesto.
Ning\'un conjunto de features adicionales puede reducir la mezcla
a una sola componente, porque la incertidumbre proviene de factores
no observables al momento de la entrada (realizaci\'on futura del
ruido $\sigma\xi$, \'ordenes institucionales no visibles, eventos
ex\'ogenos).
\end{proposition}

\subsection{Fracaso de la Optimizaci\'on por M\'inimos Cuadrados}

Un modelo de ML t\'ipico minimiza el error cuadr\'atico medio:
\begin{equation}
\hat{f} = \arg\min_f \sum_{i=1}^{N} \left(r_i - f(\bm{x}_i)\right)^2
\label{eq:mse}
\end{equation}

Bajo la mezcla~(\ref{eq:bimodal_inherent}), la soluci\'on \'optima
de~(\ref{eq:mse}) es la esperanza condicional:
\begin{equation}
f^*(\bm{x}) = \E[r \mid \bm{x}]
  = \pi_W \cdot \E[r \mid W] + (1 - \pi_W) \cdot \E[r \mid L]
\label{eq:ml_optimal}
\end{equation}

\begin{theorem}[El ML Aprende el Promedio de la Mezcla]
\label{thm:ml_fails}
Si la bimodalidad es irreducible (Proposici\'on~\ref{prop:bimodal_irreducible})
y las componentes tienen medias de magnitud similar con signos opuestos
($|\E[r|W]| \approx |\E[r|L]|$), entonces:
\begin{equation}
f^*(\bm{x}) \approx 0 \quad \forall\, \bm{x}
\end{equation}
El modelo aprende ``no hay se\~nal'' --- no porque la se\~nal no
exista, sino porque est\'a promediando dos poblaciones que se cancelan.
\end{theorem}

\begin{proof}
Con $\pi_W \approx 0.6$, $\E[r|W] \approx +2\%$ y
$\E[r|L] \approx -3\%$ (valores t\'ipicos del sistema $\af$):
$f^*(\bm{x}) = 0.6 \times 2\% + 0.4 \times (-3\%) = 0\%$.
El gradiente durante entrenamiento recibe contribuciones de signo
opuesto de winners y losers, resultando en pesos que convergen
a cero. A\~nadir m\'as features o m\'as complejidad al modelo no
resuelve el problema: si la bimodalidad es intr\'inseca al mercado,
no existe funci\'on $f(\bm{x})$ que la elimine.
\end{proof}

\begin{tcolorbox}[redbox, title=Diagn\'ostico: Por Qu\'e el ML No Aprende]
El problema fundamental es que el ML formula el trading como un
problema de \textbf{predicci\'on}: encontrar $f(\bm{x}) \to r$.
Pero bajo bimodalidad irreducible, la funci\'on $f^*$ \'optima
produce $r \approx 0$. Los losers ``contaminan'' el gradiente:
cada loser empuja los pesos en la direcci\'on opuesta a los winners.
El modelo queda paralizado en el promedio de la mezcla.

A\~nadir m\'as features, capas, o complejidad al modelo no resuelve
el problema --- lo empeora (overfitting a ruido). La soluci\'on no
es predecir mejor, sino reformular el problema.
\end{tcolorbox}


\subsection{La Reformulaci\'on: Clasificaci\'on Secuencial Post-Entrada}

El enfoque $\af$ + AEPS reformula el problema de dos maneras
fundamentales:

\begin{enumerate}[leftmargin=1.5em]
\item \textbf{No intenta predecir el retorno antes de entrar.}
  Acepta que $\pi_W \approx 0.6$ y que el 40\% de los trades van a
  fallar. La se\~nal solo identifica las \textit{condiciones} para
  la bifurcaci\'on, no predice cu\'al desenlace ocurrir\'a.

\item \textbf{Clasifica despu\'es de entrar, cuando hay m\'as
  informaci\'on.} Una vez abierto el trade, la trayectoria
  $(MFE(t), MAE(t))$ revela progresivamente a cu\'al poblaci\'on
  pertenece. El ML pregunta ``\textit{?`va a ser winner?}'' antes de
  entrar. AEPS pregunta ``\textit{?`es winner?}'' mientras el trade
  est\'a abierto. La segunda pregunta tiene respuesta porque la
  trayectoria contiene informaci\'on que los features pre-entrada
  no tienen.
\end{enumerate}

\begin{theorem}[Superioridad de la Clasificaci\'on Secuencial]
\label{thm:seq_superior}
Sea $I_{\text{pre}} = I(\bm{x}; O)$ la informaci\'on mutua entre
features pre-entrada y outcome, y sea
$I_{\text{post}}(t) = I(\text{MFE}(t); O)$ la informaci\'on mutua
entre MFE observado al tiempo $t$ y outcome. Entonces:
\begin{equation}
I_{\text{post}}(t) > I_{\text{pre}} \quad \text{para } t > t_{\text{divergencia}}
\end{equation}
donde $t_{\text{divergencia}}$ es el momento en que las trayectorias
de winners y losers empiezan a separarse.
\end{theorem}

En el sistema $\af$, $t_{\text{divergencia}} \approx 10$--$30$
minutos (tiempo medio al MFE de winners: 27 minutos). Despu\'es de
ese punto, observar el MFE proporciona m\'as informaci\'on para
clasificar el trade que \textit{cualquier} combinaci\'on de features
evaluada antes de la entrada.

Esta reformulaci\'on es an\'aloga al cambio de paradigma en
an\'alisis secuencial propuesto por Wald \cite{wald1945sequential}:
en vez de recolectar toda la evidencia y decidir al final (como el
test de Neyman-Pearson), se decide tan pronto como la evidencia
acumulada es suficiente. AEPS es un SPRT no-param\'etrico
\cite{tartakovsky2014sequential} donde el ``dato'' que se acumula
es la trayectoria MFE del trade.


% ═══════════════════════════════════════════════════════════════
\section{Condiciones Formales para que AEPS Funcione}
\label{sec:conditions}
% ═══════════════════════════════════════════════════════════════

Esta secci\'on establece las condiciones necesarias y suficientes que
\textbf{cualquier} se\~nal de entrada debe cumplir para que el sistema
AEPS genere edge positivo. Son independientes de la se\~nal $\af$
espec\'ifica: aplican a cualquier se\~nal (RSI, medias m\'oviles,
breakouts, modelos de ML, etc.) que se quiera combinar con AEPS.

\begin{tcolorbox}[keybox, title=Las Tres Condiciones de la Se\~nal]
\begin{enumerate}[leftmargin=1.5em]
\item \textbf{Separabilidad}:
  $\Delta = \KL(f_W^{\text{MFE}} \| f_L^{\text{MFE}}) > 0$

\item \textbf{Direcci\'on correcta}:
  $\E[\text{MFE} \mid W] > \E[\text{MFE} \mid L]$

\item \textbf{Separabilidad temprana}:
  la divergencia de MFE ocurre dentro de la ventana temporal
  $\tau_{\text{abort}}$ del clasificador
\end{enumerate}
Si las tres se cumplen, AEPS puede generar edge.
Si falla cualquiera, AEPS es in\'util independientemente de cu\'an
sofisticada sea la se\~nal.
\end{tcolorbox}


\subsection{Condici\'on 1: Separabilidad ($\Delta > 0$)}

\begin{definition}[Separabilidad]
La se\~nal produce separabilidad si y solo si las distribuciones de
MFE de winners y losers son estad\'isticamente distintas:
\begin{equation}
\Delta = \KL(f_W^{\text{MFE}} \| f_L^{\text{MFE}}) > 0
\end{equation}
\end{definition}

Sin separabilidad, el clasificador AEPS es trivial: TPR $=$ FPR
para todo umbral $\theta$ (Lema~\ref{lem:trivial}). El early abort
cortar\'ia la misma proporci\'on de winners que de losers, y las
mejoras en protecci\'on se cancelar\'ian con la destrucci\'on de
winners.

\textbf{Verificaci\'on emp\'irica}: $\Delta = 0.250$ nats con
$p = 0.000315$ (Mann-Whitney). La condici\'on se cumple con alta
significancia.

\textbf{?`Qu\'e se\~nales la cumplen?} Cualquier se\~nal que genere
dos clusters distinguibles en MFE. Las se\~nales de reversi\'on a
la media y breakouts tienden a cumplirla. Las se\~nales de tendencia
pura t\'ipicamente no, porque los winners se desarrollan lentamente
y su MFE temprano es indistinguible del de los losers.


\subsection{Condici\'on 2: Direcci\'on Correcta ($\E[\text{MFE}|W] > \E[\text{MFE}|L]$)}

\begin{definition}[Direcci\'on]
Los winners deben tener MFE mayor que los losers en promedio:
\begin{equation}
\E[\text{MFE} \mid W] > \E[\text{MFE} \mid L]
\end{equation}
\end{definition}

Esta condici\'on parece obvia pero no lo es. Existe un escenario
donde falla: se\~nales cuyos winners se desarrollan tard\'iamente.
Imaginemos una se\~nal donde los winners hacen un MFE bajo durante
2 horas y despu\'es explotan. En la ventana del abort (1.08h), su
MFE es similar al de los losers. AEPS los abortar\'ia como ``trades
muertos'' --- clasificar\'ia en la direcci\'on equivocada.

\textbf{Verificaci\'on emp\'irica}: $\E[\text{MFE}|W] = 3.74\%
\gg 1.81\% = \E[\text{MFE}|L]$. La condici\'on se cumple con
margen amplio.

\textbf{Conexi\'on f\'isica}: la bifurcaci\'on $\af$ garantiza
esta condici\'on por dise\~no. Cuando $\kappa$ invierte signo,
la cascada de Langevin es r\'apida e intensa (MFE alto). Cuando
$F_{\text{ext}}$ domina, el precio sigue subiendo y el MFE del
short se estanca. La f\'isica produce la asimetr\'ia.


\subsection{Condici\'on 3: Separabilidad Temprana}

\begin{definition}[Temporalidad]
La divergencia entre $f_W^{\text{MFE}}$ y $f_L^{\text{MFE}}$ debe
manifestarse dentro de la ventana temporal del clasificador:
\begin{equation}
\exists\, t^* < \tau_{\text{abort}}: \quad
\KL\!\left(f_W^{\text{MFE}(t^*)} \,\|\, f_L^{\text{MFE}(t^*)}\right) > 0
\end{equation}
donde MFE$(t^*)$ es el MFE m\'aximo alcanzado hasta el tiempo $t^*$.
\end{definition}

Esta es la condici\'on m\'as restrictiva y la que descarta a la
mayor\'ia de las se\~nales de tendencia. Si el clasificador (abort)
opera en una ventana de 1.08 horas, la informaci\'on para
distinguir winners de losers debe estar disponible dentro de esa
ventana. Si los winners tardan 4 horas en revelarse, el abort los
matar\'ia antes de que se distingan.

\textbf{Verificaci\'on emp\'irica}: el tiempo medio al MFE de
winners es 27 minutos ($\ll$ 1.08h). La separaci\'on es temprana.

\textbf{Tipolog\'ia de se\~nales por temporalidad}:
\begin{itemize}[leftmargin=1.5em]
\item \textbf{Se\~nales r\'apidas} (cumplen condici\'on 3):
  bifurcaci\'on $\af$, reversi\'on a la media (RSI2+SMA200),
  breakouts con volumen. Los desenlaces se resuelven en minutos
  a horas.
\item \textbf{Se\~nales lentas} (no cumplen condici\'on 3):
  SuperTrend, cruces de medias m\'oviles, sistemas trend-following.
  Los winners necesitan d\'ias o semanas para desarrollarse.
\end{itemize}


\subsection{Universalidad del Framework}

\begin{remark}[Generalidad del Framework AEPS]
Las tres condiciones son independientes de la se\~nal $\af$
espec\'ifica. Cualquier se\~nal que las cumpla puede usar AEPS
para generar edge. La bifurcaci\'on $\af$ las cumple por la
f\'isica del instrumento (funding rate como spring, cascadas de
liquidaci\'on), pero una se\~nal de mean reversion estad\'istica
que tambi\'en las cumpla podr\'ia usar el mismo framework AEPS
sin modificaciones.

Lo que distingue a $\af$ de se\~nales puramente estad\'isticas
(RSI, MACD, etc.) es la \textbf{robustez a cambios de r\'egimen}.
Una se\~nal estad\'istica mide desviaciones de un patr\'on hist\'orico
que puede cambiar. La se\~nal $\af$ mide \textbf{condiciones
estructurales} del instrumento --- funding rate como costo real,
OI como apalancamiento acumulado, taker ratio como flujo
direccional --- que son invariantes mientras exista el mecanismo
de futuros perpetuos con funding rate.
\end{remark}


% ═══════════════════════════════════════════════════════════════
\section{AEPS como Test Secuencial No-Param\'etrico}
\label{sec:sprt}
% ═══════════════════════════════════════════════════════════════

\subsection{An\'alogo del SPRT de Wald}

El SPRT cl\'asico \cite{wald1945sequential} decide secuencialmente
entre $H_0$ y $H_1$ acumulando el log-likelihood ratio:
\begin{equation}
\Lambda_n = \sum_{i=1}^{n} \log \frac{f_1(x_i)}{f_0(x_i)}
\end{equation}
y decide $H_1$ si $\Lambda_n \geq A$, $H_0$ si $\Lambda_n \leq B$,
y contin\'ua si $B < \Lambda_n < A$. Wald y Wolfowitz demostraron
que este test es \'optimo: minimiza el n\'umero esperado de
observaciones para niveles dados de error tipo I y II.

AEPS implementa un an\'alogo no-param\'etrico donde las ``fronteras''
$A$ y $B$ son umbrales calibrados por percentiles emp\'iricos, y
las ``observaciones'' son evaluaciones de la trayectoria
$(MFE(t), MAE(t), t)$ cada 5 segundos.

\begin{definition}[Clasificador AEPS]
Sea $\hat{C}: \bm{X}(t) \to \{W, L\}$ el clasificador:
\begin{itemize}[leftmargin=1.5em]
\item $\hat{C} = L$ (abort) si MFE$(t) < \theta_{\text{abort}}$
  despu\'es de $\tau_{\text{abort}}$ horas
\item $\hat{C} = W$ (capturar) si MFE$(t) \geq \theta_{\text{ptp}}$
\end{itemize}
Las tasas de error son:
\begin{align}
\text{TPR}(\theta) &= \Prob(\text{MFE} < \theta \mid L)
  = F_L^{\text{MFE}}(\theta) \\
\text{FPR}(\theta) &= \Prob(\text{MFE} < \theta \mid W)
  = F_W^{\text{MFE}}(\theta)
\end{align}
\end{definition}


\subsection{Mecanismos de Salida como Regiones de Decisi\'on}

Cada mecanismo cubre una zona del espacio (MFE, $t$, PnL):

\begin{table}[H]
\centering
\caption{Mecanismos de salida y su rol en la clasificaci\'on.}
\label{tab:exit_zones}
\begin{tabularx}{\textwidth}{lXcc}
\toprule
\textbf{Mecanismo} & \textbf{Rol en la clasificaci\'on} & \textbf{Valor} & \textbf{PnL medio} \\
\midrule
Early abort & Clasificador: decide $H_L$ (trade muerto) & 0.62\%, 1.08h & $-0.48\%$ \\
Stop loss & Fuera de la bimodalidad: protecci\'on extrema & $-5.30\%$ & $-4.96\%$ \\
Partial TP & Confirma $H_W$: asegura 33\% & 1.14\% & --- \\
Profit lock & Red entre PTP y trailing & 0.57\% & $+0.40\%$ \\
Trailing stop & Optimiza captura en $H_W$ confirmado & 2.0\%, 18\% cb & $+1.50\%$ \\
Pump capture & Captura \'optima con informaci\'on de exhaustion & variable & $+3.00\%$ \\
\bottomrule
\end{tabularx}
\end{table}

El stop loss opera \textbf{fuera} de la bimodalidad --- es protecci\'on
contra escenarios donde el modelo entero falla (la se\~nal se equivoc\'o
completamente). Los dem\'as mecanismos operan \textbf{dentro} de la
bimodalidad: el early abort clasifica, el partial TP confirma, el
profit lock protege la zona intermedia, el trailing captura, y el
pump capture maximiza.


\subsection{Correcci\'on de Censura}

\subsubsection{Censura por Early Abort (Losers)}

El early abort introduce censura por la derecha
\cite{kaplan1958nonparametric} en MAE y $t_{\text{MFE}}$: un trade
abortado no alcanz\'o su MAE real. La soluci\'on es usar solo losers
naturales $\mathcal{L}_N$ para calibrar SL y abort hours.

\subsubsection{Censura por Profit Lock (Winners)}

An\'alogamente, el profit lock censura winners: un trade cerrado por
profit lock tiene ETD/MFE artificialmente alto (0.724 vs 0.376 en
non-PL). La correcci\'on excluye profit lock del c\'alculo de
trailing callback:
\begin{equation}
\text{trailing\_callback} = P_{25}\!\left(\frac{\text{ETD}}{\text{MFE}}\right)_{\mathcal{W} \setminus \mathcal{W}_{\text{PL}}}
\end{equation}


% ═══════════════════════════════════════════════════════════════
\section{Teorema de Super-Aditividad}
\label{sec:superadd}
% ═══════════════════════════════════════════════════════════════

\subsection{Lemas Preparatorios}

\begin{lemma}[El Clasificador es Trivial sin Separabilidad]
\label{lem:trivial}
Si $\Delta = 0$, entonces $\forall\, \theta$:
$\mathrm{TPR}(\theta) = \mathrm{FPR}(\theta)$.
\end{lemma}

\begin{proof}
$\Delta = 0 \Rightarrow f_W^{\text{MFE}} = f_L^{\text{MFE}}$ c.s.
$\Rightarrow F_W(\theta) = F_L(\theta)$ $\forall\,\theta$.
\end{proof}

\begin{lemma}[Separabilidad Permite Clasificaci\'on Asim\'etrica]
\label{lem:asymmetric}
Si $\Delta > 0$ y $\E[\text{MFE} \mid W] > \E[\text{MFE} \mid L]$,
entonces existe $\theta^*$ tal que
$\mathrm{TPR}(\theta^*) > \mathrm{FPR}(\theta^*)$.
\end{lemma}

\begin{proof}
Dominancia estoc\'astica de primer orden: $F_L(\theta) \geq F_W(\theta)$
$\forall\,\theta$, con desigualdad estricta en al menos un punto
(por $\Delta > 0$).
\end{proof}


\subsection{Teorema Principal}

\begin{theorem}[Super-Aditividad de Se\~nal + AEPS]
\label{thm:superadd}
Sea $\mathcal{S}$ una se\~nal que cumple las tres condiciones de la
Secci\'on~\ref{sec:conditions}. Sea $\mathcal{A}$ el sistema AEPS.
Definimos:
\begin{align}
V_{\mathcal{S}} &= \E[\text{PnL} \mid \mathcal{S}, \text{exits fijos}] \\
V_{\mathcal{A}} &= \E[\text{PnL} \mid \text{entrada random}, \mathcal{A}] \\
V_{\mathcal{S}+\mathcal{A}} &= \E[\text{PnL} \mid \mathcal{S}, \mathcal{A}]
\end{align}
Entonces:
\begin{equation}
\boxed{V_{\mathcal{S}+\mathcal{A}} > V_{\mathcal{S}} + V_{\mathcal{A}}}
\end{equation}
\end{theorem}

\begin{proof}
\textbf{Paso 1}: $V_{\mathcal{S}} \leq 0$. Sin AEPS, losers corren
hasta SL. El ratio $\rho_{\text{fijo}} < (1-\pi_W)/\pi_W$. Dato
emp\'irico: $V_{\mathcal{S}} \approx -0.05\%$.

\textbf{Paso 2}: $V_{\mathcal{A}} = 0$. Con entrada random,
$\pi_W = 0.5$ y $\Delta = 0$. Por Lema~\ref{lem:trivial},
TPR $=$ FPR. La ganancia por abortar losers se cancela con la
p\'erdida por abortar winners.

\textbf{Paso 3}: $V_{\mathcal{S}+\mathcal{A}} > 0$. Con $\Delta > 0$,
por Lema~\ref{lem:asymmetric}, TPR $>$ FPR. AEPS aborta m\'as losers
que winners y optimiza captura. Emp\'iricamente:
$V_{\mathcal{S}+\mathcal{A}} \approx +0.45\%$.

Sinergia: $0.45\% - (-0.05\% + 0\%) = +0.50\% > 0$. $\square$
\end{proof}


\subsection{Ecuaci\'on Fundamental del Edge}

\begin{theorem}[Descomposici\'on del Edge]
\label{thm:edge}
\begin{equation}
\E[\text{PnL}] = \pi_W \cdot \E[\text{PnL} \mid W]
  + (1-\pi_W) \cdot \E[\text{PnL} \mid L]
\end{equation}
Rentabilidad requiere:
$\rho = |\E[\text{PnL}|W]| / |\E[\text{PnL}|L]| > (1-\pi_W)/\pi_W$.
\end{theorem}

Valores post-convergencia: $\pi_W = 0.615$, $\rho = 1.809$,
threshold $= 0.625$. Margen sobre breakeven: $+189\%$.
Profit factor: 3.53.


% ═══════════════════════════════════════════════════════════════
\section{Convergencia del Sistema}
\label{sec:convergence}
% ═══════════════════════════════════════════════════════════════

\subsection{Fases de Convergencia}

\begin{definition}[Fases]
\begin{enumerate}[leftmargin=1.5em]
\item \textbf{Cold start} ($n < 15$): par\'ametros est\'aticos.
  Objetivo: supervivencia.
\item \textbf{Warmup blend} ($15 \leq n < 25$):
  $\theta_{\text{eff}} = (1-\alpha)\theta_0 + \alpha\hat\theta$,
  $\alpha = (n-15)/10$.
\item \textbf{Adaptativo} ($n \geq 25$): 100\% emp\'irico.
\end{enumerate}
\end{definition}

\begin{proposition}[Condici\'on de Convergencia]
El sistema converge si: (1) $\Delta > 0$ persiste durante cold start;
(2) el capital sobrevive la fase 1. La bimodalidad es propiedad de la
se\~nal, no del AEPS --- existe desde el trade 1.
\end{proposition}


\subsection{Sizing Adaptativo por Confianza}

\begin{equation}
f_{\text{capital}} = f_{\text{base}} \times \min\!\left(1, \frac{n}{N_{\text{window}}}\right)
\end{equation}

Limita el da\~no durante cold start: trade 5 $\to$ sizing 1\%;
trade 15 $\to$ 3\%; trade 25+ $\to$ 5\%.


\subsection{Meta-Prior para Re-Convergencia}

\begin{definition}[Meta-Prior por Ratios Normalizados]
Memoria larga ($N = 100$--$200$) almacena ratios par\'ametro/ATR:
\begin{equation}
r_{\text{SL}} = \text{SL}_{\text{\'opt}} / \overline{\text{ATR}}, \quad
r_{\text{PTP}} = \text{PTP}_{\text{\'opt}} / \overline{\text{ATR}}, \quad \ldots
\end{equation}
Ante reset: $\text{SL}_{\text{nuevo}} = r_{\text{SL}} \times \text{ATR}_{\text{actual}}$.
Los ratios son m\'as estables que los valores absolutos porque est\'an
normalizados por volatilidad.
\end{definition}


% ═══════════════════════════════════════════════════════════════
\section{Procedimiento de Dry Run Pre-Operativo}
\label{sec:dryrun}
% ═══════════════════════════════════════════════════════════════

\subsection{Rompiendo la Circularidad}

El dry run resuelve el cold start midiendo la distribuci\'on
\textit{pura} de la se\~nal sin contaminaci\'on de exits.

\begin{definition}[Protocolo de Dry Run]
\begin{enumerate}[leftmargin=1.5em]
\item Paper mode con hold fijo de 2--3 horas. Sin SL, sin trailing,
  sin abort. Solo timeout.
\item Ejecutar $N_{\text{dry}} = 30$--$40$ trades registrando
  MFE, MAE, $t_{\text{MFE}}$.
\item Calcular $\Delta$, Mann-Whitney, E-ratio.
\end{enumerate}
\end{definition}

\subsection{Test de Admisi\'on}

\begin{definition}[Criterio de Admisi\'on]
El sistema se habilita si y solo si:
\begin{equation}
\boxed{\Delta > 0.15 \quad \wedge \quad p_{\text{MW}} < 0.10
  \quad \wedge \quad \E[\text{MFE} \mid W] > \E[\text{MFE} \mid L]}
\end{equation}
Si falla, la se\~nal no produce bimodalidad explotable. No operar.
\end{definition}

\subsection{Par\'ametros Iniciales Derivados}

Si pasa el test, los par\'ametros iniciales se derivan de la
distribuci\'on de dry run con las mismas f\'ormulas que AEPS:

\begin{table}[H]
\centering
\caption{Par\'ametros iniciales derivados del dry run.}
\begin{tabularx}{\textwidth}{lXl}
\toprule
\textbf{Par\'ametro} & \textbf{F\'ormula} & \textbf{Fallback (ATR)} \\
\midrule
\param{stop\_loss} & $P_{75}(|\text{MAE}|_L) + 0.5\%$ & $1.5 \times \text{ATR}$ \\
\param{partial\_tp} & $P_{25}(\text{MFE}_W)$ & $0.4 \times \text{ATR}$ \\
\param{trail\_activation} & $P_{50}(\text{MFE}_W)$ & $0.7 \times \text{ATR}$ \\
\param{abort\_mfe} & $P_{50}(\text{MFE}_L)$ & $0.2 \times \text{ATR}$ \\
\param{abort\_hours} & $P_{75}(t_{\text{MFE},L}) / 3600$ & 2.0h \\
\param{trail\_callback} & $P_{25}(\text{ETD}/\text{MFE})_{W \setminus \text{PL}}$ & 0.30 \\
\bottomrule
\end{tabularx}
\end{table}


% ═══════════════════════════════════════════════════════════════
\section{Monitor de Salud del Sistema}
\label{sec:health}
% ═══════════════════════════════════════════════════════════════

\subsection{Sistema de Tres Alarmas}

\subsubsection{Alarma R\'apida: Equity Drawdown}

\begin{equation}
\text{DD}(t) = \frac{E_{\max} - E(t)}{E_{\max}} > 0.10
\;\Rightarrow\; \text{PAUSAR}
\end{equation}

Con sizing 5\% y SL 5.3\%, $\sim$4 SLs consecutivos tocan 10\%.
No es mala suerte --- es cambio estructural.

\subsubsection{Alarma Media: Conteo de Stop Losses}

\begin{align}
n_{\text{SL}}(8) \geq 3 &\;\Rightarrow\; \text{sizing 50\%} \\
n_{\text{SL}}(12) \geq 5 &\;\Rightarrow\; \text{PAUSAR}
\end{align}

Si AEPS funciona, los losers mueren por abort ($-0.5\%$), no por
SL ($-5.3\%$). M\'as SLs indica p\'erdida de discriminaci\'on.

\subsubsection{Alarma Lenta: $\KL$ Rolling}

Cada 5 trades, calcular $\Delta$ sobre los \'ultimos 20:
\begin{align}
\Delta > 0.25 &\;\Rightarrow\; \text{operar normal} \\
0.10 < \Delta \leq 0.25 &\;\Rightarrow\; \text{sizing 50\%} \\
\Delta \leq 0.10 &\;\Rightarrow\; \text{bimodalidad perdida, PAUSAR}
\end{align}

\'Unica alarma que distingue mala suerte de p\'erdida estructural
de bimodalidad.

\subsection{Flujo de Control}

\begin{figure}[H]
\centering
\begin{tikzpicture}[node distance=1cm, font=\small,
  proc/.style={rectangle, rounded corners, minimum width=3.5cm, minimum height=1cm, align=center, draw=cdark, fill=clight},
  dec/.style={diamond, aspect=2.5, minimum width=2.5cm, align=center, draw=cdark, fill=corange!15, font=\small},
  arr/.style={thick,->,>=stealth}
]
  \node[proc, fill=cblue!15] (dry) {\textbf{DRY MODE}\\30--40 trades\\hold fijo, sin exits};
  \node[dec, below=1.2cm of dry] (test) {$\Delta > 0.15$?\\$p < 0.10$?};
  \node[proc, fill=cred!15, right=2.5cm of test] (stop) {\textbf{NO OPERAR}\\Se\~nal no sirve};
  \node[proc, fill=cgreen!15, below=1.2cm of test] (live) {\textbf{LIVE MODE}\\AEPS pre-calibrado\\3 alarmas activas};
  \node[dec, below=1.2cm of live] (health) {Alarma\\activa?};
  \node[proc, below left=1.2cm and 0.5cm of health, fill=cgreen!10] (ok) {Operar normal\\Sizing 100\%};
  \node[proc, below=1.2cm of health, fill=corange!15] (warn) {Reducir sizing\\50\%};
  \node[proc, below right=1.2cm and 0.5cm of health, fill=cred!15] (pause) {PAUSAR\\Volver a dry};
  \draw[arr] (dry) -- (test);
  \draw[arr] (test) -- node[above]{No} (stop);
  \draw[arr] (test) -- node[left]{S\'i} (live);
  \draw[arr] (live) -- (health);
  \draw[arr] (health) -- node[left,font=\scriptsize]{No alarma} (ok);
  \draw[arr] (health) -- node[left,font=\scriptsize]{Media} (warn);
  \draw[arr] (health) -- node[right,font=\scriptsize]{Cr\'itica} (pause);
  \draw[arr, dashed, cgray] (pause.east) -- ++(2,0) |- (dry.east);
  \draw[arr, dashed, cgray] (ok.west) -- ++(-1.5,0) |- (live.west);
\end{tikzpicture}
\caption{Flujo de control completo. Pausa cr\'itica $\to$ dry mode $\to$ re-evaluar $\Delta$.}
\label{fig:control_flow}
\end{figure}


% ═══════════════════════════════════════════════════════════════
\section{Bimodalidad por R\'egimen vs Bimodalidad por Estructura}
\label{sec:regime_vs_structure}
% ═══════════════════════════════════════════════════════════════

\subsection{La Distinci\'on Fundamental}

No toda bimodalidad es igual. Proponemos una distinci\'on cr\'itica
entre dos tipos de bimodalidad que, hasta donde sabemos, no ha sido
formalizada previamente:

\begin{definition}[Bimodalidad por R\'egimen]
\label{def:regime_bimodal}
Una se\~nal exhibe bimodalidad por r\'egimen cuando la separaci\'on
entre $f_W$ y $f_L$ depende del estado direccional del mercado.
En un mercado bajista, shorts random producen $\Delta > 0$ porque
la tendencia genera winners sistem\'aticamente. Cuando el r\'egimen
cambia (mercado alcista o lateral), $\Delta \to 0$ y la bimodalidad
desaparece.
\end{definition}

\begin{definition}[Bimodalidad por Estructura]
\label{def:struct_bimodal}
Una se\~nal exhibe bimodalidad por estructura cuando la separaci\'on
entre $f_W$ y $f_L$ proviene de las \textbf{condiciones del instrumento}
al momento de la entrada, no del r\'egimen general del mercado.
La se\~nal selecciona oportunidades transversalmente sobre un universo
de instrumentos, entrando solo en aquellos que exhiben la condici\'on
estructural independientemente de su tendencia individual.
\end{definition}

La diferencia operativa es la \textbf{robustez a cambios de r\'egimen}:

\begin{itemize}[leftmargin=1.5em]
\item Una se\~nal con bimodalidad por r\'egimen funciona mientras el
  mercado va en la direcci\'on favorable. Cuando cambia, la $\Delta$
  colapsa.
\item Una se\~nal con bimodalidad por estructura funciona en cualquier
  r\'egimen porque siempre est\'a seleccionando los instrumentos que
  \textit{ahora mismo} exhiben la condici\'on, independientemente de
  la tendencia general.
\end{itemize}


\subsection{La Se\~nal $\af$ como Generador de Bimodalidad Estructural}

La se\~nal $\af$ opera transversalmente sobre un universo de $500+$
tokens de futuros perpetuos. No apuesta a que un token espec\'ifico
va a caer --- identifica cu\'ales tokens \textit{en este momento}
tienen las condiciones estructurales para una bifurcaci\'on:

\begin{itemize}[leftmargin=1.5em]
\item Funding rate alto = costo real que pagan los longs
\item OI creciente = apalancamiento acumulado real
\item Taker ratio alto = flujo direccional de \'ordenes
\item Volumen elevado = actividad an\'omala
\item Exhaustion = se\~nales de agotamiento del impulso
\end{itemize}

Estos observables son \textbf{invariantes de r\'egimen}: un funding
de $0.05\%$ significa lo mismo en un mercado alcista que bajista.
Lo que cambia es \textit{cu\'antos} tokens exhiben la condici\'on
(m\'as en mercados alcistas euf\'oricos, menos en bajistas), pero
cuando la condici\'on existe, la din\'amica de bifurcaci\'on es la misma.

Esto explica por qu\'e los 46 trades reales del sistema muestran
bimodalidad significativa ($p = 0.000315$) a pesar de ser sobre
tokens completamente distintos (RDNT, PLAY, CYS, BULLA, etc.) en
diferentes momentos. Cada token est\'a en un r\'egimen diferente,
pero todos comparten la misma condici\'on estructural al momento
de entrada: spring $\af$ tensionado.


\subsection{Evidencia Experimental: El Contraste DOGE}

Para ilustrar la diferencia, se realiz\'o un experimento comparando
se\~nales cl\'asicas sobre un token individual (DOGEUSDT) durante
62 d\'ias.

\subsubsection{Protocolo}

Se tomaron 1500 velas de 1 hora de DOGEUSDT (62 d\'ias) y 1000PEPEUSDT
(62 d\'ias). Para cada token se generaron m\'ultiples se\~nales de
entrada (RSI, Bollinger Bands, volumen, taker ratio, pumps) y un
\textbf{control random} (entradas al azar). Para cada entrada se
midi\'o MFE y MAE con hold fijo de 3 horas, sin ning\'un mecanismo
de exit. Se calcul\'o $\Delta$, Mann-Whitney, y el score $Q$ de
bimodalidad.

\subsubsection{Resultados}

\begin{table}[H]
\centering
\caption{Bimodalidad de se\~nales sobre DOGEUSDT ($-25.6\%$ en el per\'iodo).}
\label{tab:doge_exp}
\begin{tabular}{lcccccl}
\toprule
\textbf{Se\~nal} & $\bm{N}$ & \textbf{WR} & $\bm{\Delta}$ & $\bm{p}_{\text{MW}}$ & $\bm{Q}$ & \textbf{Tipo} \\
\midrule
BB$>$1 + Vol$>$2$\times$ & 23 & 35\% & 0.557 & 0.0012 & 1.091 & R\'egimen \\
RSI14 $>$ 75 & 19 & 37\% & 0.496 & 0.0010 & 0.912 & R\'egimen \\
MR: RSI2$>$95+SMA200 & 34 & 41\% & 0.497 & 0.0001 & 0.865 & R\'egimen \\
\textbf{RANDOM (control)} & \textbf{120} & \textbf{52\%} & \textbf{0.619} & \textbf{0.0000} & \textbf{0.834} & \textbf{R\'egimen} \\
\midrule
Se\~nales $\af$ (multi-obs.) & 2--6 & --- & \multicolumn{4}{l}{Insuficientes trades por token} \\
\bottomrule
\end{tabular}
\end{table}

\begin{table}[H]
\centering
\caption{Bimodalidad de se\~nales sobre 1000PEPEUSDT ($-30.1\%$ en el per\'iodo).}
\label{tab:pepe_exp}
\begin{tabular}{lcccccl}
\toprule
\textbf{Se\~nal} & $\bm{N}$ & \textbf{WR} & $\bm{\Delta}$ & $\bm{p}_{\text{MW}}$ & $\bm{Q}$ & \textbf{Tipo} \\
\midrule
\textbf{RANDOM (control)} & \textbf{120} & \textbf{52\%} & \textbf{0.619} & \textbf{0.0000} & \textbf{0.833} & \textbf{R\'egimen} \\
BB$>$1 + Vol$>$2$\times$ & 18 & 28\% & 0.354 & 0.0266 & 0.599 & R\'egimen \\
RSI14 $>$ 75 & 15 & 33\% & 0.354 & 0.0071 & 0.554 & R\'egimen \\
\midrule
Se\~nales $\af$ (multi-obs.) & 1--2 & --- & \multicolumn{4}{l}{Insuficientes trades por token} \\
\bottomrule
\end{tabular}
\end{table}

\begin{tcolorbox}[redbox, title=Hallazgo Cr\'itico: El Random Gana]
El control random produce $Q = 0.83$ en \textbf{ambos} tokens ---
superando a la mayor\'ia de las se\~nales. Esto ocurre porque ambos
tokens cayeron $25$--$30\%$ durante el per\'iodo: \textbf{cualquier
short random es rentable en un mercado bajista}. La bimodalidad
observada no es propiedad de la se\~nal sino del r\'egimen. Cuando
el mercado cambie de direcci\'on, esta bimodalidad desaparecer\'a y
todas las se\~nales (incluidas BB, RSI, etc.) dejar\'an de funcionar
sobre estos tokens.
\end{tcolorbox}

\subsubsection{Contraste con el Sistema $\af$ Real}

\begin{table}[H]
\centering
\caption{Comparaci\'on: se\~nales single-token vs $\af$ multi-token.}
\label{tab:contrast}
\begin{tabular}{lccccl}
\toprule
\textbf{Sistema} & $\bm{N}$ & $\bm{\Delta}$ & $\bm{p}_{\text{MW}}$ &
  \textbf{Tokens} & \textbf{Tipo bimodalidad} \\
\midrule
Random en DOGE & 120 & 0.619 & 0.0000 & 1 token & R\'egimen ($-25\%$) \\
BB+Vol en DOGE & 23 & 0.557 & 0.0012 & 1 token & R\'egimen \\
RSI14$>$75 en DOGE & 19 & 0.496 & 0.0010 & 1 token & R\'egimen \\
\midrule
\textbf{$\af$ real} & \textbf{46} & \textbf{0.250} & $\bm{0.000315}$
  & \textbf{30+ tokens} & \textbf{Estructura} \\
\bottomrule
\end{tabular}
\end{table}

La $\Delta$ del sistema $\af$ real (0.250) es menor que la de las
se\~nales single-token en DOGE (0.5--0.6), pero tiene una diferencia
cualitativa fundamental: se midi\'o sobre \textbf{30+ tokens distintos}
en diferentes momentos y reg\'imenes. No es un artefacto del mercado
bajista --- es una propiedad de la se\~nal operando transversalmente.

\subsection{Cuarta Condici\'on: Robustez a R\'egimen}

El experimento motiva agregar una \textbf{cuarta condici\'on} al
framework de la Secci\'on~\ref{sec:conditions}:

\begin{tcolorbox}[keybox, title=Condici\'on 4: Robustez a R\'egimen]
Una se\~nal es explotable por AEPS de forma \textbf{sostenible}
(no solo en un r\'egimen favorable) si y solo si:
\begin{equation}
\boxed{\Delta_{\text{se\~nal}} > \Delta_{\text{random}}
  \quad \text{sobre el mismo universo de instrumentos}}
\end{equation}
donde $\Delta_{\text{random}}$ es la separabilidad de entradas
aleatorias sobre el mismo conjunto de tokens y per\'iodo.
Si $\Delta_{\text{se\~nal}} \leq \Delta_{\text{random}}$, la
bimodalidad es por r\'egimen y desaparecer\'a con el cambio de mercado.
\end{tcolorbox}

\begin{remark}[Test Operativo de Robustez]
El dry run (Secci\'on~\ref{sec:dryrun}) debe incluir un control
random: generar $N$ entradas aleatorias sobre los mismos tokens y
per\'iodo, medir su $\Delta_{\text{random}}$, y verificar que la
se\~nal lo supere. Si no lo supera, la se\~nal no tiene edge propio
--- est\'a surfeando el r\'egimen.
\end{remark}

\begin{remark}[Por Qu\'e $\af$ Cumple la Condici\'on 4]
La se\~nal $\af$ opera transversalmente sobre $500+$ tokens,
seleccionando en cada momento los que exhiben la condici\'on
estructural (funding $+$ OI $+$ pump $+$ exhaustion). Esto
diversifica autom\'aticamente entre reg\'imenes: si el mercado
general es alcista, hay m\'as tokens con la condici\'on (m\'as
oportunidades); si es bajista, hay menos (menos trades pero la
condici\'on sigue siendo v\'alida cuando aparece). La bimodalidad
no depende de la direcci\'on del mercado sino de la existencia
de exceso de leverage --- que es una condici\'on microestructural
del instrumento, no macroecon\'omica.
\end{remark}


% ═══════════════════════════════════════════════════════════════
\section{Validaci\'on Emp\'irica ($N=46$)}
\label{sec:validation}
% ═══════════════════════════════════════════════════════════════

\subsection{Evoluci\'on Temporal}

\begin{table}[H]
\centering
\caption{Evoluci\'on del sistema por per\'iodos.}
\begin{tabular}{lcccccc}
\toprule
\textbf{Per\'iodo} & $\bm{N}$ & \textbf{WR} & $\overline{\textbf{PnL}}$ &
  \textbf{PF} & $\bm{\Delta}$ & $\bm{p}_{\text{MW}}$ \\
\midrule
Trades 1--23 (pre-AEPS) & 23 & 57\% & $-0.05\%$ & $\sim$1.0 & 0.405 & 0.083 \\
Trades 24--46 (con AEPS) & 23 & 61\% & $-0.08\%$ & $\sim$0.9 & 0.551 & 0.003 \\
\'Ultimos 13 (post-conv.) & 13 & 69\% & $+0.45\%$ & 3.53 & --- & --- \\
\midrule
Total & 46 & 59\% & $-0.06\%$ & 0.99 & 0.250 & 0.000315 \\
\bottomrule
\end{tabular}
\end{table}

La se\~nal produce bimodalidad desde el inicio ($\Delta$ mejora de
0.405 a 0.551). El performance econ\'omico es breakeven hasta que
AEPS converge. Post-convergencia: PF = 3.53.

\subsection{Performance por Exit}

\begin{table}[H]
\centering
\caption{Estad\'isticas por exit reason ($N=46$).}
\begin{tabular}{lccccr}
\toprule
\textbf{Exit} & $\bm{n}$ & \textbf{WR} & $\overline{\textbf{PnL}}$ &
  \textbf{PMFE} & \textbf{Sum \$} \\
\midrule
pump\_capture & 9 & 100\% & $+3.00\%$ & 0.859 & $+\$28.57$ \\
trailing\_stop & 18 & 83\% & $+1.29\%$ & 0.227 & $+\$26.72$ \\
profit\_lock & 3 & 100\% & $+0.40\%$ & 0.277 & $+\$2.38$ \\
early\_abort & 5 & 20\% & $-0.48\%$ & --- & $-\$2.65$ \\
stop\_loss & 10 & 0\% & $-4.96\%$ & --- & $-\$53.05$ \\
oi\_abort & 1 & 0\% & $-2.42\%$ & --- & $-\$2.83$ \\
\bottomrule
\end{tabular}
\end{table}

El 88\% de las p\'erdidas provienen de SLs. La reducci\'on de SLs
via early abort ($-\$0.50$ vs $-\$5.30$) es el mecanismo principal
de generaci\'on de edge.


% ═══════════════════════════════════════════════════════════════
\section{Resumen y Arquitectura Te\'orica}
\label{sec:summary}
% ═══════════════════════════════════════════════════════════════

\begin{tcolorbox}[keybox, title=Arquitectura Te\'orica Completa]
\begin{enumerate}[leftmargin=1.5em]
\item \textbf{Bimodalidad inherente}: toda se\~nal de entrada genera
  una distribuci\'on de mezcla de resultados. El ML predictivo
  fracasa porque aprende el promedio de la mezcla ($\approx 0$).
  La soluci\'on es clasificaci\'on secuencial post-entrada.

\item \textbf{Cuatro condiciones de la se\~nal}: (1) separabilidad
  ($\Delta > 0$), (2) direcci\'on correcta ($\E[\text{MFE}|W] > \E[\text{MFE}|L]$),
  (3) separabilidad temprana (divergencia dentro de $\tau_{\text{abort}}$),
  y (4) robustez a r\'egimen ($\Delta_{\text{se\~nal}} > \Delta_{\text{random}}$).
  Las tres primeras son necesarias para que AEPS funcione; la cuarta
  es necesaria para que funcione de forma sostenible.

\item \textbf{Bimodalidad por estructura vs por r\'egimen}: las
  se\~nales cl\'asicas (RSI, BB) producen bimodalidad dependiente
  del r\'egimen de mercado. La se\~nal $\af$ produce bimodalidad
  por estructura del instrumento, operando transversalmente sobre
  $500+$ tokens. Evidencia experimental: el control random iguala o
  supera a se\~nales cl\'asicas en tokens individuales durante
  per\'iodos direccionales.

\item \textbf{Se\~nal $\af$}: cumple las cuatro condiciones por
  dise\~no f\'isico (Langevin). Mide condiciones estructurales del
  instrumento, no patrones estad\'isticos.

\item \textbf{AEPS}: clasificador secuencial no-param\'etrico
  (an\'alogo SPRT de Wald) con correcciones de censura (Kaplan-Meier)
  para early abort (losers) y profit lock (winners).

\item \textbf{Super-aditividad}: $V_{\mathcal{S}+\mathcal{A}} >
  V_{\mathcal{S}} + V_{\mathcal{A}}$. La se\~nal provee separabilidad,
  AEPS la explota. Ninguno funciona solo.

\item \textbf{Convergencia}: sizing adaptativo + dry run pre-operativo
  + meta-prior por ratios ATR.

\item \textbf{Monitor de salud}: drawdown, SL count, $\KL$ rolling.
  Detecta p\'erdida de bimodalidad antes de da\~no significativo.
\end{enumerate}
\end{tcolorbox}


% ═══════════════════════════════════════════════════════════════
\begin{thebibliography}{12}
% ═══════════════════════════════════════════════════════════════

\bibitem{epeloa2026bifurcation}
J.~Epeloa,
``Trading the $\alpha_f$ Bifurcation: Short Strategies in Perpetual Futures,''
Technical Report, Mapplics, 2026.

\bibitem{wald1945sequential}
A.~Wald,
``Sequential Tests of Statistical Hypotheses,''
\textit{The Annals of Mathematical Statistics}, vol.~16, no.~2,
pp.~117--186, 1945.

\bibitem{hamilton1989regime}
J.~D.~Hamilton,
``A New Approach to the Economic Analysis of Nonstationary Time Series
and the Business Cycle,''
\textit{Econometrica}, vol.~57, no.~2, pp.~357--384, 1989.

\bibitem{kaplan1958nonparametric}
E.~L.~Kaplan and P.~Meier,
``Nonparametric Estimation from Incomplete Observations,''
\textit{Journal of the American Statistical Association}, vol.~53,
no.~282, pp.~457--481, 1958.

\bibitem{cover2006elements}
T.~M.~Cover and J.~A.~Thomas,
\textit{Elements of Information Theory}, 2nd ed.,
Wiley-Interscience, 2006.

\bibitem{tartakovsky2014sequential}
A.~Tartakovsky, I.~Nikiforov, and M.~Basseville,
\textit{Sequential Analysis: Hypothesis Testing and Changepoint Detection},
CRC Press, 2014.

\bibitem{leung2021trailing}
T.~Leung and H.~Zhang,
``Optimal Trading with a Trailing Stop,''
\textit{Applied Mathematics \& Optimization}, vol.~83,
pp.~669--698, 2021.

\bibitem{sweeney1993mae}
J.~Sweeney,
``Maximum Adverse Excursion: Analyzing Price Fluctuations for Trading
Management,'' Wiley, 1997.
(Concepto original en \textit{Technical Analysis of Stocks \& Commodities},
1993.)

\bibitem{bayesian_sl}
``Determining Optimal Stop-Loss Thresholds via Bayesian Maximum
Drawdown Analysis,'' arXiv:1609.00869, 2016.

\bibitem{palazzi2025crypto}
R.~B.~Palazzi et al.,
``Trading Games: Beating Passive Strategies in the Bullish Crypto Market,''
\textit{Journal of Futures Markets}, 2025.

\bibitem{filippos2018tpsl}
F.~Filippos et al.,
``Take Profit and Stop Loss Trading Strategies Comparison in Combination
with an MACD Trading System,''
\textit{Journal of Risk and Financial Management}, vol.~11, no.~3, 2018.

\bibitem{siegmund1985sequential}
D.~Siegmund,
\textit{Sequential Analysis: Tests and Confidence Intervals},
Springer, 1985.

\end{thebibliography}

\end{document}
